\documentclass[12pt]{article}
\usepackage[utf8]{inputenc}
\usepackage[T1]{fontenc}
\usepackage{graphicx}
\usepackage{subcaption}
\usepackage{amsmath}
\usepackage{fancyhdr}
\usepackage{geometry}
\geometry{a4paper, margin=1in}
\usepackage{dirtytalk}
\usepackage[english]{babel}
\usepackage{csquotes}
\usepackage{hyperref}
\usepackage{listings}
\usepackage{xcolor}

\definecolor{codebg}{gray}{0.96}
\lstset{
    language=bash,
    basicstyle=\ttfamily\footnotesize,
    backgroundcolor=\color{codebg},
    frame=single, rulecolor=\color{gray!40},
    breaklines=true, columns=fullflexible,
    keepspaces=true, showstringspaces=false, upquote=true
}

\setlength{\headheight}{14.5pt}
\pagestyle{fancy}
\fancyhf{}
\lhead{Operating Systems -- Lab 2}
\rhead{Florea David, FAF-242}
\cfoot{\thepage}

\begin{document}

\begin{titlepage}
\begin{center}
    \IfFileExists{image.png}{\includegraphics[width=0.3\textwidth]{image.png} \\[0.2cm]}{}

    \textbf{MINISTRY OF EDUCATION, CULTURE AND RESEARCH \\ OF THE REPUBLIC OF MOLDOVA} \\[0.3cm]

    \textbf{Technical University of Moldova \\ Faculty of Computers, Informatics and Microelectronics \\ Department of Software and Automation Engineering} \\[2cm]

    \textbf{Florea David, FAF-242}\\[0.5cm]

    \Huge \textbf{Report} \\[0.5cm]

    \large Laboratory work nr. 2 \\[0.5cm]

    \textbf{Meet the OS You Will Build} \\[3cm]

        \begin{flushright}
        \textit{Checked by:} \\
        \textbf{Reițman Patricia}, \textit{University Assistant} \\
        Department of Software and Automation Engineering,
        \\FCIM Faculty, UTM
    \end{flushright}
    \vfill

    Chișinău -- 2026
\end{center}
\end{titlepage}

\newpage

\section*{Setup}
Track B (xv6). Ubuntu 26.04.1 LTS VM, cross-compiler \texttt{riscv64-linux-gnu-gcc}.
\begin{lstlisting}[language=bash]
cd ~/xv6-riscv 2>/dev/null || \
  git clone https://github.com/mit-pdos/xv6-riscv ~/xv6-riscv
cd ~/xv6-riscv
git log -1 --format="%h %ad %s" --date=short
ls -l
riscv64-linux-gnu-gcc --version | head -1
\end{lstlisting}
\begin{lstlisting}
Cloning into '/home/david/xv6-riscv'...
06aad25 2026-09-26 Add killzero test
-rw-rw-r-- 1 david david 1174 Oct  5 14:15 LICENSE
-rw-rw-r-- 1 david david 5744 Oct  5 14:15 Makefile
-rw-rw-r-- 1 david david 2441 Oct  5 14:15 README
drwxrwxr-x 2 david david 4096 Oct  5 14:15 kernel
drwxrwxr-x 2 david david 4096 Oct  5 14:15 mkfs
-rwxrwxr-x 1 david david 6147 Oct  5 14:15 test-xv6.py
drwxrwxr-x 2 david david 4096 Oct  5 14:15 user
riscv64-linux-gnu-gcc (Ubuntu 15.2.0-16ubuntu1) 15.2.0
\end{lstlisting}

\subsection*{First build (QEMU missing)}
The compiler was already on the VM, but the emulator wasn't. The kernel and file system image still built, and only the final step failed:
\begin{lstlisting}[language=bash]
make qemu
\end{lstlisting}
\begin{lstlisting}
/bin/sh: 1: qemu-system-riscv64: not found
riscv64-linux-gnu-gcc -march=rv64gc -g -ffile-prefix-map=/home/david/xv6-riscv=. -c -o kernel/entry.o kernel/entry.S
...
mkfs/mkfs fs.img README user/_cat user/_echo user/_forktest user/_grep user/_init ...
make: qemu-system-riscv64: No such file or directory
make: *** [Makefile:183: qemu] Error 127
\end{lstlisting}

After installing QEMU:
\begin{lstlisting}[language=bash]
qemu-system-riscv64 --version | head -1
\end{lstlisting}
\begin{lstlisting}
QEMU emulator version 10.2.1 (Debian 1:10.2.1+ds-1ubuntu3.2)
\end{lstlisting}

\section*{Part 1 -- Build and Boot}
\begin{lstlisting}[language=bash]
make qemu
\end{lstlisting}
\begin{lstlisting}
qemu-system-riscv64 -machine virt -bios none -kernel kernel/kernel -m 128M -smp 3 -nographic -global virtio-mmio.force-legacy=false -drive file=fs.img,if=none,format=raw,id=x0 -device virtio-blk-device,drive=x0,bus=virtio-mmio-bus.0

xv6 kernel is booting

hart 1 starting
hart 2 starting
init: starting sh
$ 
\end{lstlisting}
QEMU emulates a RISC-V \texttt{virt} machine with 3 cores (\texttt{-smp 3}, hence harts 1 and 2 starting next to hart 0) and 128\,MB of RAM, and uses \texttt{fs.img} as its disk. Ctrl-A, X exits the emulator.

\section*{Part 2 -- Using xv6}
\begin{lstlisting}[language=bash]
$ ls
\end{lstlisting}
\begin{lstlisting}
.              1 1 1024
..             1 1 1024
README         2 2 2441
cat            2 3 36728
echo           2 4 35592
sleep          2 5 35280
pingpong       2 6 36672
forktest       2 7 18136
grep           2 8 44184
init           2 9 35976
kill           2 10 35512
ln             2 11 35304
ls             2 12 42880
mkdir          2 13 35568
rm             2 14 35552
sh             2 15 58312
stressfs       2 16 36416
usertests      2 17 206800
grind          2 18 51824
wc             2 19 37728
zombie         2 20 34848
logstress      2 21 37560
forphan        2 22 36352
dorphan        2 23 35800
sync           2 24 34944
console        3 25 0
\end{lstlisting}
The columns are name, type (1 = directory, 2 = file, 3 = device), inode number and size in bytes.

\begin{lstlisting}[language=bash]
$ cat README
$ echo hello xv6
$ ls | grep c
$ wc README
\end{lstlisting}
\begin{lstlisting}
xv6 is a re-implementation of Dennis Ritchie's and Ken Thompson's Unix
Version 6 (v6).  xv6 loosely follows the structure and style of v6,
but is implemented for a modern RISC-V multiprocessor using ANSI C.
...
hello xv6
cat            2 3 36728
echo           2 4 35592
wc             2 19 37728
sync           2 24 34944
console        3 25 0
48 336 2441 README
\end{lstlisting}

\begin{lstlisting}[language=bash]
$ usertests -q
\end{lstlisting}
\begin{lstlisting}
usertests starting
test copyin: OK
test copyout: OK
...
test kernmem: usertrap(): unexpected scause 0xd pid=5593
...
test lazy_sbrk: OK
test partial_write: OK
test unlinkcwd: OK
ALL TESTS PASSED
\end{lstlisting}

\subsection*{Observations}
\textbf{(1) Three programs xv6 ships with.} \texttt{ls}, \texttt{cat} and \texttt{echo} (also \texttt{grep}, \texttt{wc}, \texttt{sh}, \texttt{kill}, \texttt{mkdir}, \texttt{rm}, \texttt{usertests}, ...). \texttt{ls} listed 22 programs in the root directory, plus \texttt{README} and the \texttt{console} device.

\textbf{(2) What a pipe needs.} For \texttt{ls | grep c} to work the OS must provide:
\begin{itemize}
    \item \textbf{process creation} (\texttt{fork} + \texttt{exec}): the shell starts \texttt{ls} and \texttt{grep} as two separate processes running at the same time;
    \item \textbf{inter-process communication} (the \texttt{pipe} system call plus file descriptors and \texttt{dup}): a kernel buffer connects the stdout of \texttt{ls} to the stdin of \texttt{grep}.
\end{itemize}

\textbf{(3) xv6 shell vs.\ Linux shell.} The xv6 \texttt{sh} handles the same basics as bash (running programs, \texttt{|}, \texttt{<}/\texttt{>}, \texttt{;}, \texttt{\&}), but it is a 499-line program with no variables, job control, history, tab completion or scripting.

\textbf{Interesting output:} \texttt{usertests -q} passed all tests, but along the way it printed lines like \\\texttt{usertrap(): unexpected scause 0xd}. Those come from tests such as \texttt{kernmem} and \texttt{stacktest} that deliberately make a bad memory access. They check that the kernel kills only the faulting process and keeps running.

\section*{Part 3 -- Reading the Source}
\begin{lstlisting}[language=bash]
sed -n "1,30p" user/cat.c
\end{lstlisting}
\begin{lstlisting}[language=C]
#include "kernel/types.h"
#include "kernel/fcntl.h"
#include "user/user.h"

char buf[512];

void
cat(int fd)
{
  int n;

  while ((n = read(fd, buf, sizeof(buf))) > 0) {
    if (write(1, buf, n) != n) {
      fprintf(2, "cat: write error\n");
      exit(1);
    }
  }
  if (n < 0) {
    fprintf(2, "cat: read error\n");
    exit(1);
  }
}
\end{lstlisting}

\begin{lstlisting}[language=bash]
cat user/user.h          # system-call part
\end{lstlisting}
\begin{lstlisting}[language=C]
// system calls
int fork(void);
int exit(int) __attribute__((noreturn));
int wait(int *);
int pipe(int *);
int write(int, const void *, int);
int read(int, void *, int);
int close(int);
int kill(int);
int exec(const char *, char **);
int open(const char *, int);
...
int getpid(void);
char *sys_sbrk(int, int);
int pause(int);
int uptime(void);
int sync(void);
\end{lstlisting}

\begin{lstlisting}[language=bash]
grep -n "sys_read\|sys_write" kernel/sysfile.c | head
grep -n -B1 -A12 "^sys_read" kernel/sysfile.c
\end{lstlisting}
\begin{lstlisting}[language=C]
69:sys_read(void)
83:sys_write(void)
68-uint64
69:sys_read(void)
70-{
71-  struct file *f;
72-  int n;
73-  uint64 p;
74-
75-  argaddr(1, &p);
76-  argint(2, &n);
77-  if (argfd(0, 0, &f) < 0)
78-    return -1;
79-  return fileread(f, p, n);
80-}
\end{lstlisting}

\begin{lstlisting}[language=bash]
grep -n "sys_read\|sys_write\|sys_pause\|SYS_read\|SYS_write\|SYS_pause" kernel/syscall.c kernel/syscall.h
\end{lstlisting}
\begin{lstlisting}
kernel/syscall.c:88:extern uint64 sys_read(void);
kernel/syscall.c:96:extern uint64 sys_pause(void);
kernel/syscall.c:99:extern uint64 sys_write(void);
kernel/syscall.c:115:  [SYS_read]    = sys_read,
kernel/syscall.c:123:  [SYS_pause]   = sys_pause,
kernel/syscall.c:126:  [SYS_write]   = sys_write,
kernel/syscall.h:6:#define SYS_read   5
kernel/syscall.h:14:#define SYS_pause  13
kernel/syscall.h:17:#define SYS_write  16
\end{lstlisting}

\subsection*{Observations}
\textbf{(1) System calls in \texttt{user/cat.c}.}
\begin{center}
\begin{tabular}{ll}
\hline
\textbf{call} & \textbf{asks the kernel to...} \\
\hline
\texttt{open(path, O\_RDONLY)} & open the named file and return a file descriptor \\
\texttt{read(fd, buf, n)} & copy up to \texttt{n} bytes from the file into the buffer \\
\texttt{write(1, buf, n)} & copy \texttt{n} bytes from the buffer to fd 1 (stdout) \\
\texttt{close(fd)} & release the file descriptor \\
\texttt{exit(status)} & end the process and report a status to its parent \\
\hline
\end{tabular}
\end{center}
\texttt{fprintf} is not a system call. It's a library function from \texttt{user/printf.c} that ends up calling \texttt{write}.

\textbf{(2) Where \texttt{sys\_read} is implemented.} In \textbf{\texttt{kernel/sysfile.c}, line 69} (the function body is lines 69--80; \texttt{sys\_write} is at line 83). The user-side stub generated from \texttt{user/usys.pl} puts the number \texttt{SYS\_read} = 5 in a register and executes \texttt{ecall}. The kernel's table in \texttt{kernel/syscall.c} (line 115) then maps number 5 to \texttt{sys\_read}, which fetches the arguments and calls \texttt{fileread}.

\textbf{(3) \texttt{kernel/} vs.\ \texttt{user/}.} \texttt{kernel/} is the OS itself and runs in privileged (supervisor) mode with direct access to hardware and all memory. \texttt{user/} contains ordinary programs that run unprivileged and can only get things done by asking the kernel through system calls.

\textbf{Interesting output:} the whole OS is small. \texttt{kernel/} has 44 source files with 6\,605 lines, and \texttt{user/} has 26 files with 5\,635 lines (almost two thirds of those, 3\,538, are the test suite \texttt{usertests.c}).

\section*{Part 4 -- sleep}
Current xv6 renamed the \texttt{sleep} system call to \textbf{\texttt{pause}}, because the kernel already has its own internal \texttt{sleep()} for waiting on events. \texttt{user/user.h} declares \texttt{int pause(int);}, and the kernel side is \texttt{sys\_pause} in \texttt{kernel/sysproc.c}. So my program calls \texttt{pause()} instead of the handout's \texttt{sleep()}.

\textbf{Creating \texttt{user/sleep.c}:}
\begin{lstlisting}[language=C]
#include "kernel/types.h"
#include "user/user.h"

int
main(int argc, char *argv[])
{
  if(argc != 2){
    fprintf(2, "usage: sleep <ticks>\n");
    exit(1);
  }
  pause(atoi(argv[1]));   // system call into the kernel (sleep() was renamed pause() in current xv6)
  exit(0);
}
\end{lstlisting}

\textbf{Adding \texttt{sleep} (and \texttt{pingpong}) to \texttt{UPROGS} in the \texttt{Makefile}:}
\begin{lstlisting}[language=bash]
sed -i 's|^\t\$U/_echo\\$|&\n\t$U/_sleep\\\n\t$U/_pingpong\\|' Makefile
grep -n "_echo\|_sleep\|_pingpong" Makefile | cat -A
\end{lstlisting}
\begin{lstlisting}
134:^I$U/_echo\$
135:^I$U/_sleep\$
136:^I$U/_pingpong\$
\end{lstlisting}
(\texttt{cat -A} shows the tab as \texttt{\^{}I} and the line end as \texttt{\$}, confirming the tab indentation and the trailing backslash.)

\textbf{Rebuilding:}
\begin{lstlisting}[language=bash]
make kernel/kernel fs.img
\end{lstlisting}
\begin{lstlisting}
riscv64-linux-gnu-gcc ... -c -o user/sleep.o user/sleep.c
riscv64-linux-gnu-ld -z max-page-size=4096 -T user/user.ld -o user/_sleep user/sleep.o user/ulib.o user/usys.o user/printf.o user/umalloc.o
...
mkfs/mkfs fs.img README user/_cat user/_echo user/_sleep user/_pingpong user/_forktest ...
\end{lstlisting}

\textbf{Running \texttt{sleep} in QEMU} (each output line is timestamped on the host):
\begin{lstlisting}[language=bash]
$ sleep
$ sleep abc
$ echo before
$ sleep 10
$ echo after
$ echo start; sleep 50; echo done
\end{lstlisting}
\begin{lstlisting}
[14:38:22.267] usage: sleep <ticks>
[14:38:26.186] before
[14:38:27.102] after
[14:39:07.160] start
[14:39:12.251] done
\end{lstlisting}
Without an argument the program prints its usage message and exits. \texttt{sleep 10} paused about 0.92\,s and \texttt{sleep 50} about 5.09\,s, so one xv6 timer tick is about 100\,ms. \texttt{sleep abc} returns immediately, because \texttt{atoi("abc")} is 0.

\subsection*{Optional -- pingpong}
\begin{lstlisting}[language=C]
#include "kernel/types.h"
#include "user/user.h"

int
main(int argc, char *argv[])
{
  int p2c[2], c2p[2];   // parent->child and child->parent pipes
  char b = 'x';

  if(pipe(p2c) < 0 || pipe(c2p) < 0){
    fprintf(2, "pingpong: pipe failed\n");
    exit(1);
  }

  int pid = fork();
  if(pid < 0){
    fprintf(2, "pingpong: fork failed\n");
    exit(1);
  }

  if(pid == 0){
    close(p2c[1]);
    close(c2p[0]);
    if(read(p2c[0], &b, 1) != 1){
      fprintf(2, "pingpong: child read failed\n");
      exit(1);
    }
    printf("%d: received ping\n", getpid());
    if(write(c2p[1], &b, 1) != 1){
      fprintf(2, "pingpong: child write failed\n");
      exit(1);
    }
    close(p2c[0]);
    close(c2p[1]);
    exit(0);
  }

  close(p2c[0]);
  close(c2p[1]);
  if(write(p2c[1], &b, 1) != 1){
    fprintf(2, "pingpong: parent write failed\n");
    exit(1);
  }
  if(read(c2p[0], &b, 1) != 1){
    fprintf(2, "pingpong: parent read failed\n");
    exit(1);
  }
  printf("%d: received pong\n", getpid());
  close(p2c[1]);
  close(c2p[0]);
  wait(0);
  exit(0);
}
\end{lstlisting}
\begin{lstlisting}
$ pingpong
10: received ping
9: received pong
\end{lstlisting}
The child (PID 10) blocks in \texttt{read} until the parent (PID 9) writes the byte, so \say{ping} is always printed before \say{pong}.

\break

\section*{Conclusion}
In this laboratory work I built xv6 from source with the RISC-V cross-compiler and booted it in QEMU on an emulated 3-core RISC-V machine. Inside it I used the familiar Unix tools (\texttt{ls}, \texttt{cat}, \texttt{echo}, \texttt{wc}) and a pipe, and the kernel test suite passed. Reading the source showed the split at the heart of an OS. \texttt{user/cat.c} only calls \texttt{read} and \texttt{write}; a stub turns each call into a numbered \texttt{ecall}, and \texttt{kernel/syscall.c} dispatches it to \texttt{sys\_read} in \texttt{kernel/sysfile.c}. Finally I wrote my own user program, \texttt{sleep}, which uses the system call now named \texttt{pause}, and the optional \texttt{pingpong}, which uses \texttt{fork} and two pipes. Both run correctly in the kernel I compiled myself.

\end{document}
